% !TEX root = ../monograph.tex
% Fragment for Chapter17. Root owns integration and section numbering.
\section{Exact current transport and the scalar regulator}
\label{ch17:current-suppliers}

Several ingredients of the trace transfer can be established before its
quantitative current estimates. We give them separately, since they use
the actual difference operator and the actual path action. They do not
identify a lattice coefficient by comparison with a separately chosen
continuum action.

\subsection{Difference-preserving block averages}

Let $L$ be a block ratio and $s_a$ a unit fine translation. Write
$\delta_a=s_a-1$, $\delta_a^{\rm c}=s_a^L-1$, and let
$\nabla_a=\sum_{j=1}^4c_j(s_a^j-s_a^{-j})$ be the eighth-order
centered difference of Section~\ref{ch17:current}. Write
$\nabla_a^{\rm c}=\sum_{j=1}^4c_j(s_a^{Lj}-s_a^{-Lj})$ for the
corresponding coarse difference, acting on the fine lattice.
The coefficients satisfy
$2\sum_jjc_j=1$. Define finite Laurent polynomials
\begin{align*}
 M_a&=L^{-1}\sum_{v=0}^{L-1}s_a^v,&
 G_a&=\sum_{j=1}^4c_j\sum_{v=-j}^{j-1}s_a^v,\\
 N_a&=G_a(s_a^L)M_a,&
 K_0&=\prod_aM_aN_a,\\
 K_a&=G_aM_a\prod_{c\ne a}M_cN_c,&
 K_a^{\rm e}&=N_a\prod_{c\ne a}M_cN_c.
\end{align*}
All these operators act on the fine anchor of a local function before
stationary substitution. If $\mathcal I F=F(U_{\rm stat}(V))$, coarse
translation of $\mathcal IF$ is $\mathcal I(s_a^LF)$; fine unit shifts
are not silently identified with coarse shifts.

\begin{lemma}[Exact block identities]
\label{ch17:block-identities}
There is a finite Laurent polynomial $K_a^{\rm t}$ such that
\begin{align}
 K_0\nabla_a&=L^{-1}\nabla_a^{\rm c}K_a,&
 K_0\delta_a&=L^{-1}\delta_a^{\rm c}K_a^{\rm e},&
 K_a-K_0&=\delta_a^{\rm c}K_a^{\rm t}.
 \label{ch17:block-algebra}
\end{align}
The kernels have support of radius $CL$ and absolute coefficient sums
bounded by a numerical constant independent of $L$; their masses are one
except that no mass is prescribed for $K_a^{\rm t}$.
\end{lemma}
\begin{proof}
The identities $G_a\delta_a=\nabla_a$ and
$M_a\delta_a=L^{-1}\delta_a^{\rm c}$ prove the first two assertions
by commutation of Laurent polynomials. For the third, work in the
quotient ring with $s_a^L=1$. There $M_a^2=M_a$,
$M_aG_a=G_a(1)M_a=M_a$ and $G_a(s_a^L)=G_a(1)=1$.
Thus $M_a(G_a-G_a(s_a^L)M_a)$ is zero in that quotient, which is
precisely divisibility of $K_a-K_0$ by $s_a^L-1$.
The original kernels have finite products of a fixed number of bounded
coefficient-sum factors, each of range $CL$. For the quotient, group
coefficients in each residue class modulo $L$. Its numerator has zero
sum in each class, and only a bounded number of nonzero entries in that
class. Division by $s_a^L-1$ is a partial sum along this finite sequence,
so increases its absolute sum by only a bounded factor. This also proves
the support bound for $K_a^{\rm t}$.
\end{proof}

For a fine current with $\operatorname{tr}J=\sum_a\delta_aq_a$, set
\[
 J'_{ab}=L^4\mathcal I(K_aJ_{ab}),\qquad
 q'_a=L^3\mathcal I(K_a^{\rm e}q_a)
                       +L^4\mathcal I(K_a^{\rm t}J_{aa}).
\]
Equation~\eqref{ch17:block-algebra} proves exactly
$\operatorname{tr}J'=\sum_a\delta_a^{\rm c}q'_a$.
Indeed split $K_a=K_0+\delta_a^{\rm c}K_a^{\rm t}$ and then apply
the second identity to $K_0\delta_aq_a$.
Similarly, if $J_{ab}-J_{ba}=\sum_c\delta_ct_{abc}$, the transferred
antisymmetric part is the sum of the transferred old differences and
\[
 \delta_a^{\rm c} L^4\mathcal I(K_a^{\rm t}J_{ab})
 -\delta_b^{\rm c} L^4\mathcal I(K_b^{\rm t}J_{ba}).
\]
In particular symmetry may fail pointwise after blocking, but its failure
and the trace remain exact differences. This conclusion is independent
of estimates on the current's Taylor coefficients.

\subsection{What scalar mass regularization changes}

Consider the same admissible path action used in Chapter~\ref{ch07:matching},
with its left Hessian $H$ and infinitesimal gauge insertion $d$.
At the zero field the completed vector operator is
$H+dd^*=|k|^2I+O(|k|^3)$, and the ghost operator is
$d^*d=|k|^2I+O(|k|^4)$, with the global lattice ellipticity stated there.
Define $\Gamma_\mu$ with the scalar mass $\mu^2I$ as above.
The normalized quadratic response of $\Gamma_\mu$ has the same logarithmic
coefficient as the regulated response in Chapter~\ref{ch07:matching}.
This fact concerns its coefficient integral, before any relation to the
Haar divergence of a current.

To see this directly, retain the actual vector and ghost first and second
background insertions $b,c$ and replace only their denominators by the
scalar-mass denominators. In each channel the logarithmic Hessian is
\[
 \operatorname{tr}\bigl(a_\mu(k)^{-1}c(-p,p;k)
 -a_\mu(k)^{-1}b(-p;k+p)a_\mu(k+p)^{-1}b(p;k)\bigr).
\]
The scalar mass has no background insertion. The actual local Taylor
expansion and vertex estimates of Chapter~\ref{ch07:matching} therefore
apply without the regulator-vertex terms. With $T=|k|+\mu$, the
inverse and its $r$ momentum derivatives cost $CT^{-2-r}$, for every
fixed required $r$. Differentiate $Ra=I$ repeatedly to obtain this bound;
the denominator floor is now simply $a_\mu\ge cT^2I$.
Two external derivatives of the displayed expression cost at most
$CT^{-4}$. For $\mu\le|k|\le1$, comparison with the massless continuum
vertices by the resolvent identity costs
\[
 C\bigl(|k|^{-3}+\mu^2|k|^{-6}\bigr).
\]
The first term is the actual lattice Taylor remainder, and the second
is the scalar-mass denominator difference. Both have bounded integrals
in four dimensions, uniformly in $\mu$. The region $|k|<\mu$ also
has bounded integral, since its volume is $O(\mu^4)$ and its integrand
is $O(\mu^{-4})$. Away from zero the integral is uniformly bounded.
Thus the continuum homogeneous term alone determines the logarithm.
With the twelve-plane coefficient convention of this chapter this gives
\begin{equation}
 R(\mu)=12\kappa_G\log\mu+O(1),\qquad0<\mu\le1.
 \label{ch17:scalar-log-supplier}
\end{equation}
This is the same one-loop normalization already calculated using the
actual action in Chapter~\ref{ch07:matching}. At fixed $\mu>0$ the
integrand is periodic analytic on a common smaller strip, so its torus
sums converge to this integral; no uniform rate as $\mu\downarrow0$
is required. This proof does not differentiate the bounded remainder.

\subsection{Exact integration by parts and the remaining current estimate}

For a smooth vector field $Y$ on a finite product of compact link groups,
Haar integration gives the exact identity
\begin{equation}
 \int\bigl[\operatorname{div}_{\rm Haar}Y
                  -D\mathcal A[Y]\bigr]e^{-\mathcal A}\,dU=0.
 \label{ch17:finite-haar}
\end{equation}
It follows by integrating the divergence of $e^{-\mathcal A}Y$.
With a spectator $F$ the right side becomes
$-\int DF[Y]e^{-\mathcal A}\,dU$. In particular it is zero when
$Y$ acts only on links disjoint from those of $F$. Smooth truncations
can be inserted before this identity is used; their differentiated terms
must be retained as marked remainders.

There is also an exact cancellation available for the inherited quadratic
loop term. Chapter~\ref{ch07:matching} represents its Hessian as a sum
with two lattice differences on each leg. Consequently it annihilates
constant and affine polynomial link fields. The degree-three external
momentum coefficient in its variation by the leading current pairs a
constant leg with a quadratic leg, or two affine legs; both vanish.
This proves the zero trace-functional coefficient of that quadratic
loop contribution. It does not bound the nonlinear current remainder.

The missing quantitative step is now quite specific. One needs the
actual stationary pullback of the current defect, with the gauge-covariant block average
frame correction, in the exponentially weighted order-two, -three and
-four coefficient norms. In particular the order-two pullback must
annihilate its affine--affine and affine--quadratic polynomial moment conditions
\emph{before} estimating the remainder. The usual smooth transport bound
$\nabla_f^rW_a=O(L^{-1-r})$ alone does not assert these cancellations.
After gauge removal the relevant affine field costs $L^{-2}$ and its
quadratic term $L^{-3}$; multiplication by the current rescaling $L^5$
would leave the unacceptable factors $L$ and $1$ from their two products.
Only removal of these exact terms leaves the $L^{-1}$ remainder used
in the proposed contraction. Fourth-order action matching must therefore
be combined with an explicit identity for the current's matched jets,
not merely with a bound on the action's matching error.

The group-coordinate identity needed for this comparison is proved next.
Its proof is independent of the contraction of the actual current
sequence, established in Section~\ref{ch17:quadratic-current}.

\subsection{The group-coordinate contact identity}
\label{ch17:contact-section}

Fix a finite lattice and a background $U$ in the product of link groups.
Use the invariant inner product on its link Lie algebra and the chart
$z\mapsto e^zU$ (edge by edge). Let $S$ be a smooth gauge-invariant action,
$H=D_z^2S(e^zU)|_{z=0}$ its symmetric left-coordinate Hessian, and
$d$ the infinitesimal gauge map at $U$. The displacement $Y$ is assumed
gauge covariant, so its local flow
$\Psi_s(U)=e^{sY(U)}U$ commutes with gauge transformations. It need not
be the exact time-$s$ flow; all following identities use only its
first derivative at zero.

Define the chart velocity and its linear part by
\[
 \mathcal Y_U(z)=\left.\partial_s\right|_0
    \log\bigl(\Psi_s(e^zU)\Psi_s(U)^{-1}\bigr),
 \qquad V=D_z\mathcal Y_U(0).
\]
In particular $\mathcal Y_U(0)=0$. This definition includes the
change of the left frame as the base point moves. Replacing $V$ by an
ordinary derivative of the coefficient function $Y$ would omit that
change. Let
\[
 J_U(v,z)=\left.\partial_s\partial_t\right|_{0}
       \log\bigl(\operatorname{gauge}_{e^{sv}}(e^{tz}U)U^{-1}\bigr),
 \qquad N_{i\alpha}=DS(U)[J_U(e_\alpha,e_i)].
\]
Indices $i$ run through link coordinates and $\alpha$ through gauge
coordinates, using the basis of Section~\ref{ch01:gauge} at each link and site. Differentiating gauge invariance first in $v$, then in $z$,
gives $Hd=-N$.

For $\mu>0$ at a background where the following inverses exist, write
\[
 R=(H+dd^*+\mu^2I)^{-1},\qquad
 G=(d^*d+\mu^2I)^{-1},
\]
\[
 \Gamma=\tfrac12\log\det(H+dd^*+\mu^2I)
                            -\log\det(d^*d+\mu^2I).
\]
Near the zero background the positive denominators of the preceding
subsection make these ordinary real analytic functions. Define a link
vector $Z$ by its action against an arbitrary covector $\ell$:
\begin{align}
 \ell(Z)={}&\tfrac12\sum_{i,j}R_{ij}\,
                   \ell\bigl(D_z^2\mathcal Y_U(0)[e_i,e_j]\bigr)\notag\\
 &-\sum_{i,\alpha}
       [Gd^*(V+V^*)R]_{\alpha i}\,\ell(J_U(e_\alpha,e_i)).
 \label{ch17:contact-Z}
\end{align}
All matrices and derivatives here belong to the same background and
coordinate chart.

\begin{proposition}[Exact contact identity]
\label{ch17:contact-identity}
With these definitions,
\begin{equation}
 \tfrac12\operatorname{tr}\bigl(RD^2(DS[Y])\bigr)
       -\operatorname{div}_{\rm Haar}Y
       =D\Gamma[Y]-\mu^2\operatorname{tr}(RV)+DS[Z].
 \label{ch17:contact}
\end{equation}
No stationarity of $U$ is assumed. In particular the $DS[Z]$ term
cannot be discarded before applying the relevant coefficient functional.
\end{proposition}
\begin{proof}
Let a dot denote differentiation along $\Psi_s(U)$, using the left
chart anew at that base point. The chain rule for
$S(\Psi_s(e^zU))$ gives
\[
 D^2(DS[Y])=\dot H+V^*H+HV+C,
 \qquad C_{ij}=DS[D_z^2\mathcal Y_U(0)[e_i,e_j]].
\]
This is an exact chart calculation, not a replacement of the group by
its tangent space. The differential of $\Psi_s$ in these moving charts
is $I+sV+o(s)$. Exponential charts have Haar density one at their origins,
so its Jacobian at the base point has derivative
$\operatorname{div}_{\rm Haar}Y=\operatorname{tr}V$.

Gauge equivariance of $\Psi_s$ gives $\dot d=Vd$: differentiate
$\Psi_s(\operatorname{gauge}_{e^v}U)
=\operatorname{gauge}_{e^v}\Psi_s(U)$ first in $v$ and then in $s$.
Consequently
\[
 D\Gamma[Y]=\tfrac12\operatorname{tr}(R\dot H)
       +\operatorname{tr}(dd^*RV)
       -\operatorname{tr}\bigl(dGd^*(V+V^*)\bigr).
\]
On the other hand the Hessian chain rule and
$RH=I-Rdd^*-\mu^2R$ yield
\[
 \tfrac12\operatorname{tr}(RD^2(DS[Y]))-\operatorname{tr}V
 =\tfrac12\operatorname{tr}(R\dot H)
       -\operatorname{tr}(Rdd^*V)-\mu^2\operatorname{tr}(RV)
       +\tfrac12\operatorname{tr}(RC).
\]
Finally $Hd=-N$ implies the exact resolvent identity
\[
                  Rd=dG+RNG.
\]
Subtract the two preceding displays. The leftover gauge terms are
\[
 -\operatorname{tr}\bigl((Rdd^*+dd^*R-2dGd^*)V\bigr)
       =-\operatorname{tr}\bigl(Gd^*(V+V^*)RN\bigr).
\]
Inserting the definition of $N$ and the term $\frac12\operatorname{tr}(RC)$
identifies their sum with $DS[Z]$ from \eqref{ch17:contact-Z}.
This proves \eqref{ch17:contact}.
\end{proof}

\subsection{Cancellation of the action-gradient terms}

The functional $\mathsf l$ above extracts the coefficient of $z^2t$
from the diagonal polarization pair at momenta
$p=ze_a$, $q=-ze_a+te_b$, and sums over $a\ne c$ and $b$.
We check why it kills the action variations in
\eqref{ch17:contact}. Work at fixed positive $\mu$ so all kernels
are analytic in momentum near zero. A gauge-invariant, translation- and
hypercubic-invariant scalar action $F$ has quadratic kernel
$c_F(|p|^2I-p\otimes p)+O(|p|^3)$ by the Ward calculation in
Chapter~\ref{ch07:matching}. This applies to both $S$ and $\Gamma$.
Let $Z_b$ be any anchored gauge-covariant displacement, including $Y_b$
or the displacement \eqref{ch17:contact-Z}. Simplicity of the gauge
algebra implies $Z_b(0)=0$. The linearized covariance identity says
$DZ_b(0)d_0=0$.

Use $i$ for its output link direction and $c$ for its input polarization.
Its joint momentum symbol, with output dual variable $p$ and input
variable $q$, therefore has expansion
\[
 (DZ_b)_{ic}(p,q)=\sum_j q_j A_{b,ic}^{\,j}+O((|p|+|q|)^2),
 \qquad A_{b,ic}^{\,j}=-A_{b,ij}^{\,c}.
\]
The factors $i$ from Fourier differentiation may be absorbed into $A$.
There is no term depending only on $p$: the identity obtained by
multiplying by the gradient symbol in $q$ forces it to vanish.
For $a\ne c$, the coefficient of $z^2t$ in the quadratic variation
$DF[Z_b]$ is, apart from the common factor $c_F$,
\[
       A_{b,cc}^{\,b}-2\delta_{ab}A_{b,cc}^{\,a}
                              +\delta_{bc}A_{b,ac}^{\,a}.
\]
This follows by inserting the two kinetic kernels on its two legs.
The first sum over $a\ne c$ gives
$3\sum_{b,c}A_{b,cc}^{\,b}$. The second gives minus twice that sum,
since $A_{c,cc}^{\,c}=0$. The last gives minus the same sum after
using $A_{c,ac}^{\,a}=-A_{c,aa}^{\,c}$ and exchanging $a,c$.
Thus
\begin{equation}
                  \mathsf l(DF[Z])=0.
 \label{ch17:gradient-cancellation}
\end{equation}
Higher momentum orders cannot contribute to the coefficient of degree
three. Gauge covariance of \eqref{ch17:contact-Z} follows by transforming
its chart tensors, inverse operators and contracted indices together.
The mass $\mu^2I$ is unchanged by these orthogonal gauge frame changes.
The calculation therefore applies to both action-gradient terms of the
exact contact identity. It is a finite-momentum cancellation, not a
statement that the corresponding functions vanish pointwise.

\subsection{Power counting the regulator remainder}

The remaining estimate has a transparent sufficient numerator condition.
Let $W^{(1)},W^{(0)}$ be the vector and ghost remainder operators obtained
by subtracting the leading curvature displacement from $V$; in the
notation of the current calculation these are
\[
 W_b^{(1)}=\operatorname{sym}(V_b-dM_b-D_b)+\tfrac12T_b^{(1)},
 \qquad W_b^{(0)}=\operatorname{sym}(M_bd)-\tfrac12T_b^{(0)}.
\]
Here $\chi$ is the indicator of the anchored lattice site,
$(M_bh)(x)=\chi(x)h_b(x)$,
$(D_bh)_a(x)=\delta_a\chi(x)h_b(x)$, and
$(T_b^{(r)}h)(x)=\delta_b\chi(x)h(x)$ on $r$-forms.
The full lattice shifts and chart tensors must be retained when
extracting their coefficient; their continuum symbols only identify
the leading terms that cancel.

For that leading calculation the gauge derivative is $d=-\partial$
and the curvature displacement is
$(V_bh)_a=\chi(\partial_bh_a-\partial_ah_b)$.
Then $dM_b+D_b=-\chi\partial_ah_b$ and
$\operatorname{sym}(\chi\partial_b)=-\frac12\partial_b\chi$.
Thus both $W^{(r)}$ have zero first-order principal symbol. Covariantizing
these expressions gives the same cancellation at constant background
to first field order: the skew-adjoint multiplication by
$\operatorname{ad}A_b$ has zero symmetric part, and the remaining
$\operatorname{ad}A_a h_b$ term cancels the corresponding term of $dM_b$.
Accordingly their zero-, one- and two-field symbols have orders
$2,1,0$, respectively. A lattice/current sequence having the stated
matched free jets inherits precisely these vanishings in its limit.
Weighted coefficient bounds control the Taylor remainders.

Here is the resulting estimate without suppressed differentiation orders.
Set $z=|q|+\mu$. For $n\le2$ background derivatives and $r\le3$ external
momentum derivatives the inverse estimates are
$Cz^{-2-n-r}$. The remainder symbols have bounds
$Cz^{2-n-r}$. The quadratic-field coefficient of
$\operatorname{tr}(RW^{(1)})+\operatorname{tr}(GW^{(0)})$, after the
three momentum derivatives defining $\mathsf l$, is therefore bounded
by a sum of terms
\[
 Cz^{-2-n-r}\,z^{2-(2-n)-(3-r)}=Cz^{-5}.
\]
In four dimensions
$\int_{|q|\le1}(|q|+\mu)^{-5}\,dq\le C/\mu$.
The region away from zero is bounded. Hence the regulator factor
$\mu^2$ gives a remainder at most $C\mu$.
For approximate matched jets, first fix $\mu\ge\delta>0$; their
errors then tend to zero uniformly on that regulator interval.
Only afterwards send $\delta$ to zero. This is the order required
by the logarithmic mean-value argument.

\subsection{The corrected coefficient and order of limits}
\label{ch17:regulator-completion}

Retain the lattice correction in $\mathsf h$. For the same unit vector
$u\in\gLie$, set
\[
 \mathsf c(F)=\tfrac12\operatorname{Re}\sum_{a\ne c}[z^2]
 \widehat{F_c}_{cu,cu}(ze_a,-ze_a),
 \qquad \mathsf h=\mathsf l-\mathsf c.
\]
The quadratic variation $DF[Z]$ in the preceding calculation starts at
momentum degree three. Thus $\mathsf c(DF[Z])=0$, in addition to
\eqref{ch17:gradient-cancellation}. Both functionals therefore annihilate
the action-gradient terms of \eqref{ch17:contact}.

Suppose now that the actual current sequence has the following
properties at the fixed output level: uniform exponentially weighted
order-two through order-four coefficient bounds; convergence of the
needed low polynomial jets to the curvature current; the exact trace
and antisymmetry difference identities; and defects tending to zero
through order four in those weighted norms in
\[
                 DS[Y_b]\equiv_4\sum_a\nabla_aJ_{ab}-\mathcal E_b.
\]
The notation $\equiv_4$ refers to equality of field derivatives through
order four at the zero background. It suffices because the contact
expression uses two internal and two external field derivatives.
At fixed $\mu\ge\delta>0$, contracted defect errors tend to zero.
Translation covariance makes contraction by the inverse $R$ commute
with translating the anchor. The $\mathsf l$ coefficient of the
contracted current divergence selects $a=b$ at its sole linear
 total-momentum derivative. The sum is the contracted trace of $J$,
which has a further total-momentum zero because it is a forward
difference. Its coefficient is zero. For $\mathsf c$, the total
momentum is zero from the outset and the divergence vanishes directly.
The exact contact identity consequently yields
\[
 \mathfrak h(Y)=\mu^2\mathsf h\bigl(\operatorname{tr}(RV_b)\bigr)+o(1),
\]
uniformly on every fixed positive regulator interval. All operators here
are functions of the background before the two field derivatives are
extracted.

The definitions of $W^{(1)},W^{(0)}$ give an exact trace decomposition.
Using the same link and site bases induced by Section~\ref{ch01:gauge}, apply $Rd=dG+RNG$ to obtain
\begin{align*}
 \operatorname{tr}(RV_b)
  ={}&\operatorname{tr}\bigl(R(D_b-\tfrac12T_b^{(1)})\bigr)
      +\tfrac12\operatorname{tr}(GT_b^{(0)})+DS[Z_b']
      +\mathrm{Rem}_b,\\
 \mathrm{Rem}_b={}&\operatorname{tr}(RW_b^{(1)})
                                   +\operatorname{tr}(GW_b^{(0)}),\\
 \ell(Z_b')={}&\sum_{i,\alpha}(GM_bR)_{\alpha i}
                                       \ell(J_U(e_\alpha,e_i)).
\end{align*}
For example, the part containing $dM_b$ becomes
$\operatorname{tr}(GM_bd)+\operatorname{tr}(RNGM_b)$;
its symmetric ghost part is precisely
$\frac12\operatorname{tr}(GT_b^{(0)})+\operatorname{tr}(GW_b^{(0)})$,
and its remaining term is $DS[Z_b']$. This also fixes the sign of
the gauge correction.

For the three displayed principal terms the Fourier transform of
$\delta_b\chi$ supplies $e^{-it}-1=-it+O(t^2)$ at total momentum
$te_b$. With the factor $i$ in $\mathsf l$, the $D_b$ terms sum to
$\mathsf q(\operatorname{tr}R)$: the selected derivative has $a=b$
and the component trace then sums its diagonal entries. Each of the
four $T_b^{(1)}$ terms gives $\mathsf q(\operatorname{tr}R)$, and each
$T_b^{(0)}$ term gives $\mathsf q(\operatorname{tr}G)$.
The volume normalization is exactly the translation average of the
single anchor. Their total is therefore
\[
           -\mathsf q(\operatorname{tr}R)
                              +2\mathsf q(\operatorname{tr}G).
\]
Their $\mathsf c$ coefficients are zero, since $\delta_b\chi$ has
zero symbol at total momentum zero. Multiplication by $\mu^2$ gives
$-\mu\partial_\mu\mathsf q(\Gamma)$ by differentiation of the two
determinants. Thus, with the lattice correction fully retained,
\begin{equation}
 \mathfrak h(Y)=\mathcal D_\Lambda(\mu)
                  +\mu^2\mathsf h(\mathrm{Rem})+o(1).
 \label{ch17:corrected-regulated-ward}
\end{equation}
The action-gradient term $DS[Z_b']$ is annihilated by both parts of
$\mathsf h$ as proved above.

The $\mathsf l$ remainder was bounded by $C\mu$ after multiplication
by $\mu^2$. The correction $\mathsf c$ uses only two external momentum
derivatives, so the same product rule costs $Cz^{-4}$. Its integral
is at most $C(1+|\log\mu|)$ in four dimensions, and
$\mu^2(1+|\log\mu|)\le C\mu$ for $0<\mu\le1$.
Consequently \eqref{ch17:corrected-regulated-ward} proves the regulator
estimate, including the anchored correction, from the stated current
hypotheses.

Finally combine it with \eqref{ch17:scalar-log-supplier}. On a fixed
interval $[u-A,u]$ in the variable $\log\mu$, the bounded error in
$R(\mu)-12\kappa_G\log\mu$ implies that some regulator has
$|\mathcal D_\Lambda(\mu)+12\kappa_G|\le(2C+o(1))/A$ after the
period is large. The remainder is at most $Ce^u+o(1)$ on that interval.
First take the current and period limits there, next $u\to-\infty$,
and finally $A\to\infty$. Since $\mathfrak h(Y)$ is independent of
$\mu$, this proves
\[
                         \mathfrak h(Y)\longrightarrow-12\kappa_G.
\]
The group-coordinate contact identity, coefficient cancellation and
regulator estimate are now established. Their application still requires
the actual transferred current sequence to satisfy the weighted bounds
and matched-jet convergence stated at the start of this subsection;
those conditions are not consequences of the contact algebra.
